% Boyer Reinterpretation
\documentclass[aps,prd,twocolumn,superscriptaddress,nofootinbib]{revtex4-2}

\usepackage{amsmath,amssymb,amsfonts}
\usepackage{hyperref}
\usepackage{booktabs}

\begin{document}

\title{Boyer's Casimir Repulsion as Euler Product Truncation:\\
Number-Theoretic Reinterpretation and\\
Three Routes to Experimental Realization}

\author{Wright Brothers Research Program}
\affiliation{Independent Research, 2026}

\date{\today}

\begin{abstract}
Boyer (1974) predicted that the Casimir force between a perfectly
electrically conducting (PEC) plate and a perfectly magnetically
conducting (PMC) plate is repulsive, in contrast to the attractive
PEC--PEC case.  This prediction remains unverified after fifty
years, primarily because no PMC material exists at the relevant
(optical--UV) frequencies.  We offer a number-theoretic
reinterpretation: the PEC--PMC (Dirichlet--Neumann) mode sum
$1+3+5+7+\cdots=\zeta_{\neg 2}(-1)=+1/12$ is precisely the
Riemann zeta function with the $p{=}2$ Euler factor removed,
while the PEC--PEC sum $1+2+3+4+\cdots=\zeta(-1)=-1/12$ is the
full zeta function.  The sign flip is thus the arithmetic operation
of \emph{prime muting}: removing the contribution of the smallest
prime from the vacuum mode sum.  This reinterpretation clarifies why
PMC is not required---any physical configuration selecting odd-integer
modes suffices.  We identify three such configurations that bypass
the PMC obstacle: (A)~nanoscale open-ended structures providing
Neumann boundary conditions without magnetic materials;
(B)~$\lambda/4$ superconducting resonators in circuit QED, where
Lamb shift differences between $\lambda/4$ and $\lambda/2$
resonators probe the vacuum energy sign; and (C)~Si-MEMS devices
with DN-type geometries, extending the platform of Tang \textit{et
al.}\ (2017).  All three are realizable with existing technology.
A positive result would constitute the first experimental
verification of Boyer's fifty-year-old prediction and the first
physical realization of Euler product truncation.
\end{abstract}

\maketitle

%===================================================================
\section{Introduction: A Fifty-Year Gap}
%===================================================================

In 1974, Boyer showed that the Casimir energy between a perfect
electric conductor (PEC, Dirichlet boundary condition for the
electric field) and a perfect magnetic conductor (PMC, Neumann
boundary condition) is \emph{positive}---opposite in sign to the
standard PEC--PEC result~\cite{Boyer1974}.  Kenneth \textit{et al.}\
provided a rigorous proof~\cite{Kenneth2002}, and Hertzberg
\textit{et al.}\ confirmed it in the piston
geometry~\cite{Hertzberg2005}.

Despite this theoretical clarity, \textbf{no experiment has ever
measured the PEC--PMC Casimir force}.  The obstacle is material:
PMC requires magnetic permeability $\mu\to\infty$ at the optical
and ultraviolet frequencies that dominate the Casimir energy.
No natural material satisfies this, and metamaterial approaches
fail due to Kramers--Kronig constraints on
broadband magnetic response~\cite{Rosa2008}.

Several experiments have claimed ``Casimir repulsion,'' but none
realizes Boyer's DN sign reversal (Table~\ref{tab:prior}).

\begin{table}[b]
\caption{Prior ``Casimir repulsion'' experiments and why none
constitutes a DN vacuum energy sign reversal.}
\label{tab:prior}
\begin{ruledtabular}
\begin{tabular}{p{1.7cm}p{2.5cm}p{3.5cm}}
Experiment & Mechanism & Why not DN repulsion \\
\hline
Munday \textit{et al.}\ (2009, \textit{Nature})~\cite{Munday2009}
& Dielectric contrast: gold--bromobenzene--silica.
$\varepsilon_1>\varepsilon_\mathrm{fluid}>\varepsilon_2$ (DLP condition).
& Mode sum unchanged ($\sum n^{-s}$, all modes present).
Repulsion from material contrast, not boundary-condition
change. No modes selectively removed. \\
Tang \textit{et al.}\ (2017, \textit{Nature Photon.})~\cite{Tang2017}
& Geometric interplay: T-shaped Si nanostructures
slide past each other. Net force vector changes sign
at certain lateral displacements.
& All pairwise Casimir interactions remain \emph{attractive}.
``Repulsion'' is the vector sum of multiple attractions
in different directions. Mode sum unchanged.
No modes removed. \\
Boyer sphere (1968)~\cite{Boyer1968}
& Conducting spherical shell.
Casimir self-energy is positive (outward pressure).
& Geometry-dependent sign (sphere vs.\ plates), but mode
sum is the full spectrum of the sphere.
No prime muting, no DN boundary condition. \\
Metamaterial proposals~\cite{Rosa2008,Leonhardt2007}
& Left-handed or chiral metamaterials
engineered to approximate PMC response.
& Kramers--Kronig: narrowband $\mu(\omega)$ resonances
are smoothed out at imaginary frequencies.
Casimir force integral sees no selectivity.
Broadband PMC is physically impossible for passive media. \\
\end{tabular}
\end{ruledtabular}
\end{table}

In every case, the underlying vacuum mode structure is
$\sum_{n=1}^{\infty}n^{-s}$: no modes are selectively removed,
and the Euler product remains intact.  These experiments
demonstrate force-direction control or material-contrast effects,
\emph{not} vacuum energy sign reversal via mode selection.

\textbf{The DN Casimir force---where the mode sum is genuinely
truncated to $\sum_{n\;\mathrm{odd}}n^{-s}=\zeta_{\neg 2}(s)$---has
never been experimentally realized.}

%===================================================================
\section{Number-Theoretic Reinterpretation}
%===================================================================

The Riemann zeta function admits an Euler product over primes:
\begin{equation}\label{eq:euler}
  \zeta(s)=\sum_{n=1}^{\infty}n^{-s}
  =\prod_{p\;\mathrm{prime}}\frac{1}{1-p^{-s}}.
\end{equation}
Removing the $p{=}2$ factor yields
\begin{equation}\label{eq:zeta2}
  \zeta_{\neg 2}(s)\equiv(1{-}2^{-s})\,\zeta(s)
  =\sum_{\substack{n=1\\n\;\mathrm{odd}}}^{\infty}n^{-s}.
\end{equation}
The physically relevant case is three-dimensional ($s=-3$),
where the Casimir energy density between parallel plates involves
$\sum n^3$:
\begin{align}
  \zeta(-3)&=1^3{+}2^3{+}3^3{+}4^3{+}\cdots=+\tfrac{1}{120},\label{eq:z3}\\
  \zeta_{\neg 2}(-3)&=1^3{+}3^3{+}5^3{+}7^3{+}\cdots=-\tfrac{7}{120}.\label{eq:z3neg}
\end{align}
The sign flip $+1/120\to-7/120$ is the
\emph{defining property} of prime-2 muting in 3D.
(In 1D, the analogous flip is
$\zeta(-1)=-1/12\to\zeta_{\neg 2}(-1)=+1/12$.)

The standard 3D Casimir energy density is
$\rho_{DD}=-\pi^2\hbar c/(720\,a^4)$, where $720=6/\zeta(-3)$.
Boyer's result for PEC--PMC is:
\begin{itemize}
\item PEC--PEC (DD): all modes
$\;\Rightarrow\;\rho\propto -\zeta(-3)<0$ (attractive).
\item PEC--PMC (DN): odd modes only
$\;\Rightarrow\;\rho\propto -\zeta_{\neg 2}(-3)=+7/120>0$ (repulsive).
\end{itemize}
The DN energy density is positive and 7 times the DD magnitude.
\textbf{Boyer computed $\zeta_{\neg 2}(-3)$ without knowing it.}
His electromagnetic boundary-condition calculation and the
number-theoretic Euler product truncation are the same
mathematical operation, viewed from different angles.

This reinterpretation has a crucial practical consequence:
\emph{PMC is not needed}.  Any physical configuration whose
mode spectrum consists of odd integers realizes $\zeta_{\neg 2}$.
The task is not to find a magnetic material, but to find a
geometry that selects odd modes.

%===================================================================
\section{The $\lambda/4$ Realization}
%===================================================================

A $\lambda/4$ resonator---one end shorted (Dirichlet), the other
open (Neumann)---has mode frequencies
\begin{equation}
  f_n=(2n{-}1)\,\frac{c}{4L},\quad n=1,2,3,\ldots
\end{equation}
These are the odd integers $1,3,5,7,\ldots$ in units of $c/(4L)$.
In 1D, the vacuum energy involves $\zeta(-1)$:
\begin{equation}
  E_{\lambda/4}=\frac{\pi\hbar c}{4L}\,\zeta_{\neg 2}(-1)
  =+\frac{\pi\hbar c}{48L}>0.
\end{equation}
Compare the $\lambda/2$ resonator (both ends shorted or both open):
\begin{equation}
  E_{\lambda/2}=\frac{\pi\hbar c}{2L}\,\zeta(-1)
  =-\frac{\pi\hbar c}{24L}<0.
\end{equation}
In 3D between parallel plates (area $A$, separation $a$), the
Casimir energy densities are:
\begin{equation}
  \rho_{DD}=-\frac{\pi^2\hbar c}{720\,a^4},\quad
  \rho_{DN}=+\frac{7\pi^2\hbar c}{720\,a^4},
\end{equation}
with $\rho_{DN}/\rho_{DD}=-7$, directly reflecting
$\zeta_{\neg 2}(-3)/\zeta(-3)=-7$.
\textbf{Opposite signs.  No PMC material.  No metamaterial.
No Kramers--Kronig obstruction.}  The open end of a transmission
line provides a broadband, dispersion-free Neumann boundary
condition for TEM modes.

The obstacle is scale: transmission-line $\lambda/4$ resonators
are macroscopic ($\sim$cm), where the Casimir force
$\sim 10^{-21}$\,N is unmeasurable.  We propose three routes
around this obstacle.

%===================================================================
\section{Three Routes to Experimental Realization}
%===================================================================

\subsection{Route A: Nanoscale DN Structures}

At the nanoscale ($a\sim 100$\,nm), Casimir forces are measurable
($\sim$pN).  The challenge is to create DN boundary conditions at
this scale.

A Neumann condition for the electromagnetic field corresponds
physically to an open aperture, a slot, or a free edge.  Candidate
structures include:
\begin{itemize}
\item Open-ended metallic nanowaveguides (slot arrays in thin metal
films).
\item Nanoscale coaxial structures with open terminations,
fabricable by template-based
electrodeposition~\cite{Nanocoax2007}.
\item Photonic crystal surfaces engineered to exhibit
high-impedance (PMC-like) behavior over a broad bandwidth.
\end{itemize}
The Casimir force between a PEC plate and such a nano-DN surface
can be measured with standard AFM techniques.

\subsection{Route B: Circuit QED ($\lambda/4$ vs.\ $\lambda/2$)}

In circuit QED, a transmon qubit coupled to a superconducting
resonator experiences a Lamb shift
\begin{equation}\label{eq:lamb}
  \delta f=\sum_n\frac{g_n^2}{f_q-f_n},
\end{equation}
where $f_n$ are the resonator mode frequencies and $g_n$ the
coupling strengths.

A $\lambda/4$ resonator (modes $1,3,5,\ldots$) and a $\lambda/2$
resonator (modes $1,2,3,\ldots$) of the same fundamental frequency
will produce \emph{different} Lamb shifts, because the mode sums
in~\eqref{eq:lamb} differ.  The difference
\begin{equation}
  \Delta(\delta f)=\delta f_{\lambda/2}-\delta f_{\lambda/4}
  =\sum_{n\;\mathrm{even}}\frac{g_n^2}{f_q-f_n}
\end{equation}
isolates the even-mode vacuum fluctuation contribution.

Both $\lambda/4$ and $\lambda/2$ resonators are standard
components in superconducting circuit laboratories.  The
measurement requires a dilution refrigerator ($\sim$10\,mK) and
a transmon qubit---existing infrastructure at dozens of
laboratories worldwide.

\subsection{Route C: MEMS with DN Geometry}

Tang \textit{et al.}~\cite{Tang2017} demonstrated on-chip Casimir
force measurement between silicon nanostructures using integrated
comb actuators and vibrating-beam force sensors.  Their T-shaped
geometry produced non-monotonic (but not genuinely DN) forces.

We propose adapting their platform with a DN-type geometry:
one surface is a flat PEC (gold-coated silicon), the other is a
periodic array of open-ended slots or channels whose mode
structure approximates DN boundary conditions.  The force is
measured as a function of separation $a$ using the same
comb-actuator technology.

The key measurement is the scaling exponent $\alpha$ in
$F\propto a^{-\alpha}$.  Standard PEC--PEC gives $\alpha=5$
(attractive, well verified).  PEC--DN should give $\alpha=5$
with \emph{opposite sign} (repulsive, never measured).

%===================================================================
\section{Predictions}
%===================================================================

\begin{enumerate}
\item \textbf{Sign}: The Casimir force in a DN configuration is
repulsive.  This is Boyer's prediction, now reinterpreted as
Euler product truncation $\zeta\to\zeta_{\neg 2}$.

\item \textbf{Magnitude}: In 3D, the ratio of DN to DD Casimir
energy densities is $\zeta_{\neg 2}(-3)/\zeta(-3)=-7$.  The DN
repulsive force is 7 times the magnitude of the DD attractive
force at the same separation.

\item \textbf{Scaling}: Standard QED predicts $F\propto a^{-5}$
for both DD and DN.  Any deviation from $\alpha=5$ in the DN
case would indicate physics beyond the standard boundary-effect
framework.
\end{enumerate}

%===================================================================
\section{Discussion}
%===================================================================

\subsection{Why This Hasn't Been Done}

The PEC--PMC Casimir repulsion has remained
unverified for fifty years because the problem was framed as
a \emph{materials} problem (``find a PMC'').  The number-theoretic
reinterpretation reframes it as a \emph{mode selection} problem
(``select odd modes''), which can be solved geometrically.

The $\lambda/4$ solution---simply opening one end of a
resonator---is arguably the simplest realization of Boyer's
prediction.  It has not been proposed in this context because
(i)~$\lambda/4$ resonators are macroscopic where Casimir forces
are negligible, and (ii)~the connection between DN modes and
Euler product truncation was not previously recognized.

\subsection{Relation to Broader Program}

The Euler product $\zeta(s)=\prod_p(1{-}p^{-s})^{-1}$
decomposes the vacuum into prime-indexed channels.  Removing the
$p{=}2$ channel ($\zeta\to\zeta_{\neg 2}$) is the simplest
nontrivial modification.  The sign flip
$\zeta(-1)\to\zeta_{\neg 2}(-1)$ ($-1/12\to+1/12$) and
$\zeta(-3)\to\zeta_{\neg 2}(-3)$ ($+1/120\to-7/120$) suggest
that prime muting could, in principle, reverse the sign of
the cosmological vacuum energy---a necessary condition for
exotic spacetime geometries~\cite{Alcubierre1994,Visser1995}.
Experimental verification of Boyer's sign flip in any of the
three proposed routes would be the first step toward testing
this broader hypothesis.

\subsection{What This Paper Does Not Claim}

We do not claim that prime muting is physically realizable
beyond the DN boundary condition.  The connection between
DN modes and the $p{=}2$ Euler factor is a mathematical
identity, not a physical mechanism.  Whether the arithmetic
structure of the Euler product has deeper physical significance
is an open question that experiment must address.

%===================================================================
\section{Conclusion}
%===================================================================

Boyer's 1974 prediction that PEC--PMC Casimir force is repulsive
can be reinterpreted as the statement that removing the $p{=}2$
Euler factor from the Riemann zeta function reverses the sign of
the regularized vacuum energy.  This reinterpretation shifts the
experimental challenge from materials (finding a PMC) to geometry
(selecting odd modes).  We have identified three experimentally
viable routes---nanoscale DN structures, circuit QED $\lambda/4$
comparison, and MEMS with DN geometry---each realizable with
existing technology.  A positive result would simultaneously
verify a fifty-year-old theoretical prediction and provide the
first physical realization of Euler product truncation in quantum
field theory.

\begin{thebibliography}{99}
\bibitem{Boyer1974} T.~H.~Boyer, Phys.\ Rev.\ A \textbf{9}, 2078 (1974).
\bibitem{Kenneth2002} O.~Kenneth, I.~Klich, A.~Mann, and M.~Revzen, Phys.\ Rev.\ Lett.\ \textbf{89}, 033001 (2002).
\bibitem{Hertzberg2005} M.~P.~Hertzberg, R.~L.~Jaffe, R.~Kardar, and A.~Scardicchio, Phys.\ Rev.\ Lett.\ \textbf{95}, 250402 (2005).
\bibitem{Rosa2008} F.~S.~S.~Rosa, D.~A.~R.~Dalvit, and P.~W.~Milonni, Phys.\ Rev.\ Lett.\ \textbf{100}, 183602 (2008).
\bibitem{Munday2009} J.~N.~Munday, F.~Capasso, and V.~A.~Parsegian, Nature \textbf{457}, 170 (2009).
\bibitem{Tang2017} L.~Tang \textit{et al.}, Nature Photon.\ \textbf{11}, 97 (2017).
\bibitem{Nanocoax2007} F.~Rizzi \textit{et al.}, in \textit{Nanocoaxial structures} (various, 2007--2015).
\bibitem{Boyer1968} T.~H.~Boyer, Phys.\ Rev.\ \textbf{174}, 1764 (1968).
\bibitem{Leonhardt2007} U.~Leonhardt and T.~G.~Philbin, New J.\ Phys.\ \textbf{9}, 254 (2007).
\bibitem{Alcubierre1994} M.~Alcubierre, Class.\ Quant.\ Grav.\ \textbf{11}, L73 (1994).
\bibitem{Visser1995} M.~Visser, \textit{Lorentzian Wormholes} (AIP, 1995).
\end{thebibliography}

\end{document}
